\documentclass{article}
\usepackage[T1]{fontenc}
\usepackage{lmodern}
\usepackage{graphicx}
\usepackage{amsmath,amssymb}
\usepackage[a4paper,margin=2cm]{geometry}
\usepackage[absolute,overlay]{textpos}
\setlength{\TPHorizModule}{1mm}
\setlength{\TPVertModule}{1mm}
\usepackage[indonesian]{babel}
\usepackage{hyperref}
\setlength{\emergencystretch}{2em}
% unit: Halaman 17

\begin{document}

% === Halaman 17 (llm) ===

\section*{Keamanan RSA}

\begin{itemize}
    \item Keamanan algoritma $RSA$ terletak pada tingkat kesulitan dalam memfaktorkan bilangan bulat $n$ faktor-faktor prima ($p$ dan $q$), yang dalam hal ini $n = p \times q$. Semakin besar $n$, semakin sulit difaktorkan.
    \item Sekali $n$ berhasil difaktorkan menjadi $p$ dan $q$, maka
    \[ \phi(n) = (p - 1) \times (q - 1) \text{ dapat dihitung.} \]
    Selanjutnya, karena kunci enkripsi $e$ diumumkan (tidak rahasia), maka kunci dekripsi $d$ dapat dihitung dari kekongruenen
    \[ ed \equiv 1 \pmod{\phi(n)} \rightarrow d = \frac{1 + k \, \phi(n)}{e} \]
\end{itemize}

\end{document}
